\documentclass[10pt,a4paper]{amsart}
\usepackage[margin=19mm]{geometry}
\usepackage{amsmath,amssymb,mathtools}
\usepackage[scr=rsfs,cal=euler]{mathalfa}
\usepackage{xfrac}
\usepackage[unicode,hidelinks,bookmarks=false]{hyperref}
\newcommand{\ie}{i.e.\ }
\newcommand{\cf}{cf.\ }
\newcommand{\resp}{respectively\ }
\newcommand{\Subsection}{\subsection}
\providecommand{\nobreakdash}{\nobreak\hbox{-}}
\setlength{\emergencystretch}{2em}
\multlinegap0pt
\usepackage{bm}
\usepackage{bbm}
\usepackage{dsfont}
\usepackage{xspace}
\usepackage{cancel}
\usepackage{mathtools}
\usepackage{amsthm}
\makeatletter
\@ifundefined{theorem}{\newtheorem{theorem}{Theorem}}{}
\@ifundefined{lemma}{\newtheorem{lemma}[theorem]{Lemma}}{}
\makeatother
\newcommand{\N}{\mathbb{N}}
\title[Quantitative homogenization of elliptic equations with infin]{Quantitative homogenization of elliptic equations with infinitely many scales}
\author{Zhongwei Shen, Yao Xu, Jinping Zhuge}
\date{Source: arXiv:2605.02561v1 (CC BY 4.0)}
\begin{document}
\maketitle

\noindent\emph{Source and licence.} Quantitative homogenization of elliptic equations with infinitely many scales. Version arXiv:2605.02561v1 (CC BY 4.0), distributed under the Creative Commons Attribution 4.0 International licence, as recorded in the arXiv licence metadata for this exact version and shown on the abstract page https://arxiv.org/abs/2605.02561v1. Full rights notice: licenses/arxiv-2605.02561-CC-BY-4.0.txt.\par\smallskip

\noindent\textit{Editorial scope.} Counted unit: the Lemma whose source label is \texttt{int-Lip\_prop\_Hh} in the article (source0.tex lines 2524--2533, source label \texttt{int-Lip\_prop\_Hh}), delivered together with the article's own complete original proof of it (source0.tex lines 2607--2709, one \texttt{proof} environment of 103 lines). No mathematics is written, repaired or re-derived: statement and proof are the article's own text line for line. The proof is detached from the statement in the source: the article states the result at lines 2524--2533 and prints its proof at lines 2607--2709, and the proof environment's own header names the statement, so the association is the article's own. Required context retained uncounted: int-Lip\_cond\_sumbound (lines 2059--2062); int-Lip\_lem\_varepsilon (lines 2074--2079); cond\_logseparation\_m (lines 2507--2509); int-Lip\_def\_MN (lines 2516--2521); lem.sumMm (lines 2537--2542). Every definition, display and label of those lines is retained unchanged. Omitted from this package and not redistributed: 1--2058, 2063--2073, 2080--2506, 2510--2515, 2522--2523, 2534--2536, 2543--2606, 2710--3134. Nothing inside a retained excerpt is omitted or altered: every retained body is delivered exactly as the article's lines are written, so no omission receipt of any kind accompanies this package. Citations: the retained excerpts carry no citation command at all, so no bibliography entry of the article is transcribed and no reference is redistributed. Reference adaptation: none was needed and none was made; every reference inside the delivered unit resolves inside it, so no unresolved reference or citation is produced. Printed slips kept literal: no printed word, symbol, hypothesis or number of the counted statement or of its proof was altered. Layout and class: the article's own class is \texttt{amsart} with packages including mathrsfs, amsmath, amsthm, amssymb, amsfonts, bbm, esint, tikz-cd, color, hyperref, geometry, comment; the wrapper is the lane's pinned \texttt{amsart} editorial wrapper at 10pt on A4, which restates the theorem-like environments the retained text needs and adds the article's own macro definitions that the retained text needs (\texttt{\textbackslash N}), with extra wrapper packages (bm, bbm, dsfont, xspace, cancel, mathtools). The delivered numbering is the unit's own. Assets: no retained span contains a figure environment, a graphics-inclusion command, a PDF-page inclusion, a table or a TikZ picture; no image file is copied into this package, the package declares no asset dependency and no image file is redistributed with it. Licence: version arXiv:2605.02561v1 of this article is distributed under the Creative Commons Attribution 4.0 International licence, as recorded in the arXiv licence metadata for this exact version and shown on the abstract page https://arxiv.org/abs/2605.02561v1. A full rights notice is delivered at licenses/arxiv-2605.02561-CC-BY-4.0.txt. Characters: every retained excerpt is pure ASCII, so the delivered engine, XeLaTeX with its default Latin Modern text font, reports no missing character.
\begin{align}
  \label{int-Lip_cond_sumbound}
  \sum_{j=1}^{m}\frac{\varepsilon_m}{\varepsilon_j}\leq K\quad \text{for each }m \in \N_+,
\end{align}

\begin{lemma}\label{int-Lip_lem_varepsilon}
  Let $1=\varepsilon_0> \varepsilon_1>\varepsilon_2>\cdots$ be a sequence of positive numbers. Suppose \eqref{int-Lip_cond_sumbound} is fulfilled for $m\geq 1$. Then, for $1\leq k\leq m$,
  \begin{align}\label{est.em/ek}
    \frac{\varepsilon_m}{\varepsilon_{k}}\leq K(1-K^{-1})^{m-k}.
  \end{align}
\end{lemma}

\begin{align}\label{cond_logseparation_m}
  \bigg(\frac{\varepsilon_k}{\varepsilon_{k-1}}\bigg)^{\alpha\sigma}\leq c_0\Big|\log \frac{\varepsilon_{m+1}}{\varepsilon_m}\Big|^{-1} \quad\text{for any } k\geq m+2,
\end{align}

\begin{align}\label{int-Lip_def_MN}
\begin{aligned}
  M_m: & = \max \Big\{ [\widetilde{\delta}]_\alpha, \sum_{k=0}^\infty \| \nabla_{x} \widetilde{B}_\ell \|_{L^\infty} \Big\} \\
  & \le \max\bigg\{\sum_{j=0}^{m}\frac{\varepsilon_m}{\varepsilon_j}\sum_{k=j}^{\infty}\delta_k, \sum_{j=0}^m\frac{\varepsilon_m}{\varepsilon_j}\cdot\sum_{k=1}^\infty k^\alpha\delta_{k+m}\bigg\}.
  \end{aligned}
\end{align}

\begin{lemma}\label{lem.sumMm}
    Assume $[\delta]_{1+\rho} < \infty$ for some $\rho>0$ and \eqref{int-Lip_cond_sumbound}. Let $\alpha = \rho/2$ in \eqref{int-Lip_def_MN}. Then for any $T > 0$, there exists $m_0 \in \N$, depending only on $\rho, T, [\delta]_0$ and $[\delta]_{1+\rho}$, such that
    \begin{equation*}
        T \sum_{\ell = m_0}^\infty M_\ell \le 1.
    \end{equation*}
\end{lemma}

\begin{lemma}\label{int-Lip_prop_Hh}
    Suppose that $[\delta]_1<\infty$, $0<\alpha\leq 1$ and \eqref{int-Lip_cond_sumbound} holds. There exists a constant $c_0\leq 1$, depending on $d, \mu, p, [\delta]_0 $ and $[\delta]_1$, such that if the scale-separation condition \eqref{cond_logseparation_m} holds for $m\geq 0$, then for all $r\in [\varepsilon_{m+1}, \varepsilon_m]$
    \begin{align*}%\label{int_es_base}
    \begin{split}
        & H(r; u_\varepsilon)\leq C_0 \Big(\frac{r}{\varepsilon_m}\Big)^\lambda H(\varepsilon_m; u_\varepsilon)+C_1M_m\{H(\varepsilon_m; u_\varepsilon)+h(\varepsilon_m; u_\varepsilon)\},\\
        & |h(r; u_\varepsilon)-h(\varepsilon_m; u_\varepsilon)|\leq C_0H(\varepsilon_m; u_\varepsilon)+C_1M_m\{H(\varepsilon_m; u_\varepsilon)+h(\varepsilon_m; u_\varepsilon)\},
    \end{split}
    \end{align*}
    where $M_m$ is given by \eqref{int-Lip_def_MN}, $\lambda\in(0, 1)$ and $C_1\geq C_0\geq 1$. Furthermore, $\lambda$ and $C_0$ depend only on $d, \mu, p $ and $[\delta]_0$, while $C_1$ depends additionally on $\alpha$ and $[\delta]_1$. 
\end{lemma}

\begin{proof}
% We set the paramter $\alpha=\rho/2$ in all $M_m$. By Lemma \ref{int-Lip_lem_varepsilon} and Remark \ref{int-lip_remark_decay}, we calculate that for $m\in \mathbb{N}$
% \begin{align*}
%   &\quad\sum_{\ell=m}^\infty\sum_{j=0}^{\ell}\frac{\varepsilon_\ell}{\varepsilon_j}\sum_{k=j}^\infty\delta_k=\sum_{\ell=m}^\infty\sum_{k=0}^{\ell}\sum_{j=0}^{k}\frac{\varepsilon_\ell}{\varepsilon_j}\delta_k+\sum_{\ell=m}^\infty\sum_{k=\ell+1}^{\infty}\sum_{j=0}^{\ell}\frac{\varepsilon_\ell}{\varepsilon_j}\delta_k\\&\leq (K+1)\sum_{\ell=m}^\infty\sum_{k=0}^{\ell}\frac{\varepsilon_\ell}{\varepsilon_{k}}\delta_k+(K+1)\sum_{k=m+1}^\infty (k-1)\delta_k\\&=(K+1)\bigg\{\sum_{\ell=m}^\infty\varepsilon_\ell\delta_0+\sum_{k=1}^{m-1}\sum_{\ell=m}^\infty\frac{\varepsilon_\ell}{\varepsilon_{k}}\delta_k+\sum_{k=m}^{\infty}\sum_{\ell=k}^\infty\frac{\varepsilon_\ell}{\varepsilon_{k}}\delta_k\bigg\}+\frac{K+1}{m^\rho}[\delta]_{1+\rho}\\&\leq (K+1)^3\bigg\{[\delta]_0(1-K^{-1})^{m-1}+\sum_{k=1}^{m-1}(1-K^{-1})^{m-k}\delta_k+\sum_{k=m}^{\infty}\delta_k\bigg\}+\frac{K+1}{m^\rho}[\delta]_{1+\rho}\\&\leq (K+1)^3\bigg\{[\delta]_0(1-K^{-1})^{m-1}+\max\{m^{-1}, (1-K^{-1})^{m-1}\}\sum_{k=1}^{m-1}k\delta_k\\&\qquad\qquad+m^{-1}\sum_{k=m}^{\infty}k\delta_k\bigg\}+\frac{K+1}{m^\rho}[\delta]_{1+\rho}\\&\leq (K+1)^3([\delta]_0+[\delta]_1)\max\{m^{-1}, (1-K^{-1})^{m-1}\}+\frac{K+1}{m^\rho}[\delta]_{1+\rho},
% \end{align*}
% and, by using \eqref{int-Lip_cond_sumbound}
% \begin{align*}
%   &\quad\sum_{\ell=m}^\infty\sum_{j=0}^\ell\frac{\varepsilon_\ell}{\varepsilon_j}\sum_{k=1}^\infty k^\alpha\delta_{k+\ell}\leq (K+1)\sum_{\ell=m}^\infty\sum_{k=1}^\infty k^\alpha\delta_{k+\ell}\\&= (K+1)\sum_{\ell=m}^\infty\sum_{k=\ell+1}^\infty k^\alpha\delta_{k}= (K+1)\sum_{k=m+1}^\infty\sum_{\ell=m}^{k-1} k^\alpha\delta_{k}\\&\leq (K+1)\sum_{k=m+1}^\infty k^{1+\alpha}\delta_k\leq \frac{K+1}{m^{\rho-\alpha}}\sum_{k=m+1}^\infty k^{1+\rho}\delta_k\leq \frac{K+1}{m^{\rho/2}}[\delta]_{1+\rho}.
% \end{align*}
% This implies that there exists $m$ large enough, depending only on $\rho, K, C_1, [\delta]_0, [\delta]_{1+\rho}$, such that
% \begin{align}\label{int-Lip_es_M_remainder}
%   16C_1^2\sum_{\ell=m}^\infty M_\ell\leq \frac{1}{4},
% \end{align}
%   where $C_1$ is given in Lemma \ref{int-Lip_prop_Hh}. In this situation,
% \begin{align}\label{int-Lip_property_Mj}
%     64C_1^2M_\ell\leq 1\quad \text{for }\ell\geq m.
% \end{align}
Let $C_1$ be given by Lemma \ref{int-Lip_prop_Hh}. By Lemma \ref{lem.sumMm}, there exists $m_0 \in \N$ such that
\begin{align}\label{int-Lip_es_M_remainder}
    64C_1^2\sum_{\ell=m_0}^\infty M_\ell\leq 1,
\end{align}
and therefore,
\begin{align}\label{int-Lip_property_Mj}
    64C_1^2 M_\ell\leq 1\quad \text{for all }\ell \geq m_0.
\end{align}
We claim that for $\ell\geq m_0 \geq 0$
  \begin{align}\label{int-Lip_H_claim}
    \begin{split}
      H(\varepsilon_{\ell+1}) & \leq C_0^{\ell-m_0+1}\Big(\frac{\varepsilon_{\ell+1}}{\varepsilon_{m_0}}\Big)^\lambda H(\varepsilon_{m_0})\\
      & \qquad+\sum_{j=m_0}^{\ell}C_0^{\ell-j}\Big(\frac{\varepsilon_{\ell+1}}{\varepsilon_{j+1}}\Big)^\lambda C_1M_j\{H(\varepsilon_j)+h(\varepsilon_j)\}.
    \end{split}
  \end{align}
  This can be proved by induction on $\ell$. The case $\ell=m_0$ is given by Lemma \ref{int-Lip_prop_Hh}. Suppose that it is true for $\ell+1$. Then, by Lemma \ref{int-Lip_prop_Hh} and using the inductive hypothesis, we have
  \begin{align*}
    H(\varepsilon_{\ell+2})&\leq C_0\Big(\frac{\varepsilon_{\ell+2}}{\varepsilon_{\ell+1}}\Big)^\lambda H(\varepsilon_{\ell+1})+C_1M_{\ell+1}\{H(\varepsilon_{\ell+1})+h(\varepsilon_{\ell+1})\}\\
    &\leq C_0\Big(\frac{\varepsilon_{\ell+2}}{\varepsilon_{\ell+1}}\Big)^\lambda \bigg\{C_0^{\ell-m_0+1}\Big(\frac{\varepsilon_{\ell+1}}{\varepsilon_{m_0}}\Big)^\lambda H(\varepsilon_{m_0})\\
    &\qquad \qquad+\sum_{j=m_0}^\ell C_0^{\ell-j}\Big(\frac{\varepsilon_{\ell+1}}{\varepsilon_{j+1}}\Big)^\lambda C_1M_j\big\{H(\varepsilon_j)+h(\varepsilon_j)\big\}\bigg\}\\
    &\qquad +C_1M_{\ell+1}\{H(\varepsilon_{\ell+1})+h(\varepsilon_{\ell+1})\}\\&\leq C_0^{\ell-m_0+2}\Big(\frac{\varepsilon_{\ell+2}}{\varepsilon_{m_0}}\Big)^\lambda H(\varepsilon_{m_0})\\&\qquad+\sum_{j=m_0}^{\ell+1} C_0^{\ell+1-j}\Big(\frac{\varepsilon_{\ell+2}}{\varepsilon_{j+1}}\Big)^\lambda C_1M_j\{H(\varepsilon_j)+h(\varepsilon_j)\},
  \end{align*}
  which is exactly the claim.

Taking the summation of \eqref{int-Lip_H_claim} from $m_0$ to $k > m_0$, we obtain by Lemma \ref{int-Lip_lem_varepsilon} that
  \begin{align*}
    \sum_{\ell=m_0}^k H(\varepsilon_{\ell+1}) &\leq \sum_{\ell=m_0}^kC_0^{\ell-m_0+1}\Big(\frac{\varepsilon_{\ell+1}}{\varepsilon_{m_0}}\Big)^\lambda H(\varepsilon_{m_0})\\&\qquad+ \sum_{\ell=m_0}^k\sum_{j=m_0}^{\ell}C_0^{\ell-j}\Big(\frac{\varepsilon_{\ell+1}}{\varepsilon_{j+1}}\Big)^\lambda C_1M_j\{H(\varepsilon_j)+h(\varepsilon_j)\}\\&\leq 2 C_0(1-K^{-1})^{\lambda} H(\varepsilon_{m_0})\sum_{\ell={m_0}}^k[C_0(1-K^{-1})^{\lambda}]^{\ell-m_0}\\&\quad+2 \sum_{j=m_0}^kC_1M_j\{H(\varepsilon_j)+h(\varepsilon_j)\}\sum_{\ell=j}^k[C_0(1-K^{-1})^{\lambda}]^{\ell-j}\\&\leq 2H(\varepsilon_{m_0})+\sum_{\ell=m_0}^k 4C_1M_\ell\{H(\varepsilon_\ell)+h(\varepsilon_\ell)\},
  \end{align*}
  where in the last step we have required $1 < K \leq K_0$ and $K_0 \in (1,2]$ is close to $1$ such that
\begin{align}\label{int-Lip_cond_K0}
  C_0(1-K_0^{-1})^{\lambda}\le 1/2.
\end{align}

Therefore, by \eqref{int-Lip_property_Mj},
\begin{align*}
    \sum_{\ell=m_0}^{k+1}H(\varepsilon_{\ell}) & = H(\varepsilon_{m_0}) + \sum_{\ell=m_0}^{k}H(\varepsilon_{\ell+1})\\
    & \leq 3H(\varepsilon_{m_0})+\frac{1}{2} \sum_{\ell = {m_0}}^k H(\varepsilon_\ell) + \sum_{\ell = m_0}^k 4C_1M_\ell h(\varepsilon_\ell),
\end{align*}
which leads to
\begin{align}\label{int-Lip_es_H}
  \sum_{\ell = m_0}^{k+1}H(\varepsilon_{\ell})\leq 6H(\varepsilon_{m_0})+\sum_{\ell = m_0}^k8C_1M_\ell h(\varepsilon_\ell).
\end{align}

On the other hand, using the estimate of $h$ in Lemma \ref{int-Lip_prop_Hh}, we have for $k\geq m_0$,
\begin{align*}
  |h(\varepsilon_{k+1})-h(\varepsilon_{m_0})| & \leq \sum_{\ell=m_0}^k|h(\varepsilon_{\ell+1})-h(\varepsilon_\ell)|\\
  &\leq C_0\sum_{\ell=m_0}^k H(\varepsilon_\ell)+\sum_{\ell=m_0}^kC_1M_\ell\{H(\varepsilon_\ell)+h(\varepsilon_\ell)\}\\
  &\leq 2C_1\sum_{\ell=m_0}^k H(\varepsilon_\ell)+\sum_{\ell=m_0}^kC_1M_\ell h(\varepsilon_\ell),
\end{align*}
where we have used the fact $M_\ell\leq 1$ for $\ell\geq m_0$ derived from \eqref{int-Lip_property_Mj}. Substituting \eqref{int-Lip_es_H} into the inequality above and using \eqref{int-Lip_es_M_remainder}, this gives
\begin{align*}
  h(\varepsilon_{k+1})&\leq h(\varepsilon_{m_0})+ 12C_1H(\varepsilon_{m_0})+\sum_{\ell=m_0}^{k}16C_1^2M_\ell h(\varepsilon_\ell)+\sum_{\ell=m_0}^kC_1M_\ell h(\varepsilon_\ell)\\
  & \leq h(\varepsilon_{m_0})+ 12C_1H(\varepsilon_{m_0})+\frac{1}{2}\sup_{m_0 \leq \ell\leq k}h(\varepsilon_\ell).
\end{align*}
Taking the supremum over $k$, we obtain
\begin{align*}
  \sup_{m_0 \leq \ell\leq k}h(\varepsilon_{\ell})\leq h(\varepsilon_{m_0})+ 12C_1H(\varepsilon_{m_0})+\frac{1}{2}\sup_{m_0 \leq \ell\leq k}h(\varepsilon_\ell),
\end{align*}
which implies, 
\begin{align*}
  \sup_{\ell \ge m_0} h(\varepsilon_{\ell})&\leq 2\{h(\varepsilon_{m_0})+ 12C_1H(\varepsilon_{m_0})\}. 
\end{align*}
Combining this estimate with \eqref{int-Lip_es_H} and \eqref{int-Lip_es_M_remainder}, and taking $k\to \infty$, we also have
\begin{align*}
  \sum_{\ell=m_0}^{\infty} H(\varepsilon_\ell)\leq 6H(\varepsilon_{m_0})+\sum_{\ell=m_0}^{\infty}8C_1M_\ell h(\varepsilon_\ell)\leq h(\varepsilon_{m_0})+8C_1H(\varepsilon_{m_0}). 
\end{align*}
Consequently, we conclude that
\begin{align}\label{est.m0-infity}
  \sum_{\ell = m_0}^\infty H(\varepsilon_\ell) + \sup_{\ell \ge m_0} h(\varepsilon_\ell)\leq C\{h(\varepsilon_{m_0})+H(\varepsilon_{m_0})\},
\end{align}
where $C$ depends only on $C_1$.

Furthermore, using the estimates in Lemma \ref{int-Lip_prop_Hh} recursively, it holds for $1\leq \ell\leq m_0$
  \begin{align}\label{est.1-m_0}
    H(\varepsilon_{\ell})+h(\varepsilon_{\ell})\leq C\{H(\varepsilon_{\ell-1})+h(\varepsilon_{\ell-1})\}\leq\cdots\leq  C^{\ell}\{H(1)+h(1)\}, 
  \end{align}
  where $C$ depends only on $d, \mu, p, \rho, [\delta]_0$ and $[\delta]_1$. According to Lemma \ref{int-Lip_lem_varepsilon}, we know $\varepsilon_k\rightarrow 0$ as $k\rightarrow\infty$. This means that for any $r\in (0, 1]$ there exists $k\in\mathbb{N}$ such that $\varepsilon_{k+1}<r\leq \varepsilon_k$. Using Lemma \ref{int-Lip_prop_Hh} again, we have
  \begin{align}\label{est.in2scale}
    H(r)+h(r)\leq C\{H(\varepsilon_k)+h(\varepsilon_k)\}. 
  \end{align}
Combining \eqref{est.m0-infity}, \eqref{est.1-m_0} and \eqref{est.in2scale}, we finally obtain
  \begin{align*}
    H(r)+h(r)\leq C\{H(1)+h(1)\}\quad\text{for all } r\in (0,1), 
  \end{align*}
  where $C$ also depends on $m_0$. This completes the proof.
\end{proof}

\end{document}
